% chap2.tex
\part{유지보수 및 신기술 적용}
\label{part:module}

% ---------------------------------------------------------------------------- %
%                                  NEW SECTION                                 %
% ---------------------------------------------------------------------------- %

\chapter{알고리즘의 수정과 코드의 유지보수}

\section{알고리즘의 수정 제안}

이런 값들을 바꿔보면서 수정하면 됨

% ---------------------------------------------------------------------------- %
%                                  NEW SECTION                                 %
% ---------------------------------------------------------------------------- %

\section{기본값 변경}

to\_idf\_object 

% ---------------------------------------------------------------------------- %
%                                  NEW SECTION                                 %
% ---------------------------------------------------------------------------- %

\section{일반적인 코드의 유지보수}

유지보수 방법

% ---------------------------------------------------------------------------- %
%                                  NEW SECTION                                 %
% ---------------------------------------------------------------------------- %

\chapter{신기술 적용}

\section{신기술의 시뮬레이터 적용하기 위해 필요한 개념적인 과정}

지금 체계에 맞아야 함. 예를 들어서,

% ---------------------------------------------------------------------------- %
%                                  NEW SECTION                                 %
% ---------------------------------------------------------------------------- %

\section{신기술을 본 시뮬레이터로 테스트해보는 방법}

idf에만 모듈 형식으로 결합할 수 있도록 작성. 일단 idf로 내보낸 다음에 신기술 override하고 결과해석 모듈 돌리는 방법이 있음.

% ---------------------------------------------------------------------------- %
%                                  NEW SECTION                                 %
% ---------------------------------------------------------------------------- %

\section{제도적으로 제출하기 위해 필요한 것}

이런 것들을 테스트하여 제출하기 바람. 테스트 건물은 우리가 쓴 것 쓰면 좋을 듯.